% Template for post or presentation abstracts submitted to ILAS 2022
% 1. Fill in the presenter's information (name, affiliation, email
% address) and talk information (title of presentation, abstract).
% 2. Compile to make sure there are no errors.
% 3. Save the .tex file for your abstract with a distinctive name,
% ideally "FamilynameGivenname.tex".
% 4. Upload your .tex file to https://clr.ie/129557
%       
% * Please keep your LaTeX commands simple, and avoid the use of
%    user-defined macros.
% * Include details of co-authors and funding acknowledgements after the abstract.
% * Upload only the LaTeX file (with .tex extension).
% * Questions: Contact Niall Madden (Niall.Madden@NUIGalway.ie)

\documentclass[a4paper]{amsart} 
\usepackage{amssymb} 
\usepackage{hyperref}
\pagestyle{empty}
\date{} 

%% IGNORE THE NEXT 4 LINES
\makeatletter % 
\def\labelenumi{\theenumi.}
\def\theenumi{\@arabic\c@enumi}
\makeatother
%% STOP IGNORING NOW

\begin{document}

\title{The Generalized Notion of Frobenius Normal Form of Powers of Matrices}
\author{Mashael AlBaidani} %Presenting author only
\address{Prince Sattam Bin Abdulaziz University} % Affiliation of presenting author
\email{m.albaidani@psau.edu.sa} % Email address of presenting author only

\maketitle

% The abstract  ***no more than one page***
 
  Barker \cite{barker1} attempted to create a Frobenius normal form for cone nonnegative matrices. In this talk, we present conditions on a cone $\mathcal{K}$ so that any $\mathcal{K}$-nonnegative matrix admits a Frobenius normal form, and provide an example to show that not every cone admits this property. Moreover, based on the spectral characterization of those real matrices that leave a proper polyhedral cone invariant, we provide the generators for a proper polyhedral cone $\mathcal{K}$, such that $A$ is eventually $\mathcal{K}$-irreducible with a small number of extreme rays and how we can construct the Frobenius normal form with respect to cones with the aid the idea of level sets.
  



%% Coauthor details here (optional - delete if not applicable)
% and funding acknowledgment (optional - delete if not applicable)
%\emph{This is joint work with John Todd (Caltech) and Alexander Ostrowski (Basel).Supported by the National Mathematics Foundation, Grant XYZ-123.}

\begin{thebibliography}{1}
\bibitem{barker1}
George P Barker.
\newblock Stochastic matrices over cones.
\newblock {\em Linear and Multilinear Algebra}  1:279- 287,  (1974).


\end{thebibliography}


\end{document}


% Please leave the next bit alone  
%%% Local Variables: 
%%% mode: latex
%%% TeX-master: "../ILAS-2022"
%%% End:
